\chapter{رفتار زمان اجرا و کنترل محیط}

\section{راه‌اندازی}
پس از راه‌اندازی کلاک و پیرامونی‌ها، همه محرک‌های فعال‌پایین خاموش می‌شوند، گذرگاه‌های دما آزاد می‌شوند، جدول شناسه بارگذاری می‌شود و مدار انرژی به حالت \lr{SPI} قفل، بازنشانی و اجرا می‌رود. سپس \lr{DMA}، دریافت تک‌بایتی \lr{UART} و سه تایمر آغاز می‌شوند. تأخیر پنج‌ثانیه‌ای قدیمی و ارسال \lr{I2C} با اشاره‌گر نامعتبر حذف شده‌اند.

\section{کنترل رطوبت}
در حالت خودکار، پایین‌تر از ۹۸ درصد نقطه تنظیم رطوبت‌ساز روشن و از ۱۰۱ درصد آن به بالا خاموش می‌شود. بین دو مرز وضعیت قبلی حفظ می‌شود. در نسخه قدیمی خرابی حسگر دیجیتال مانع کنترل بود، درحالی‌که حسگر فعال آنالوگ است؛ وابستگی نادرست حذف شده است.

\section{کنترل سرمایش}
اگر دمای محیط بیش از نیم درجه بالاتر از نقطه تنظیم باشد، دو خروجی سرمایش روشن می‌شوند. اگر دما نیم درجه یا بیشتر پایین‌تر باشد، خاموش می‌شوند. ناحیه یک‌درجه‌ای میانی هیسترزیس از قطع‌و‌وصل سریع جلوگیری می‌کند. حفاظت حداقل زمان خاموش یا روشن کمپرسور هنوز وجود ندارد و باید پیش از اتصال مستقیم کمپرسور توسط سخت‌افزار یا نرم‌افزار جداگانه تضمین شود.

\section{کنترل گرمایش}
زیر ۹۵ درصد نقطه تنظیم، هیتر پیوسته روشن و از ۱۰۱ درصد به بالا خاموش است. در ناحیه میانی، رابطه تجربی کد قدیمی به یک پنجره صدمرحله‌ای محدود شده است. این امکان روشن‌ماندن سرمایش و جبران با هیتر را حفظ می‌کند. رابطه تجربی در نقطه تنظیم تقریباً ۱۶ درصد زمان روشن می‌دهد؛ در لبه پایین می‌تواند نزدیک صد درصد شود.

\section{حالت‌ها و ایمنی}
فرمان
\lr{AUTOFF}
کنترل خودکار را غیرفعال و همه خروجی‌های شناخته‌شده را خاموش می‌کند. فرمان
\lr{AUTONN}
آن را فعال می‌کند. فرمان نقطه تنظیم نیز کنترل را مجاز می‌سازد. فرمان‌های دستی مستقیماً خروجی را تغییر می‌دهند، ولی اگر حالت خودکار فعال بماند چرخه بعدی می‌تواند آن را بازنویسی کند؛ نرم‌افزار رایانه باید پیش از کنترل دستی، حالت خودکار را خاموش کند.

\warningbox{معنای الکتریکی خروجی‌ها از روی کد «فعال‌پایین» استنتاج شده است. پیش از نخستین پروگرام نسخه بازآرایی‌شده، حالت خاموش واقعی هر رله و \lr{SSR} با مولتی‌متر بررسی شود.}
